\section{Limitations and Future Work}
\label{sec:limitations}

\textbf{Reactive state resets on topology change.} The circuit for an entire network is
rebuilt from scratch whenever its wiring changes anywhere within that network, which resets
a capacitor's stored charge or an inductor's current to zero at that instant. The memristor
is deliberately exempted, since persistent state is its central pedagogical purpose, but a
student who extends a circuit while a capacitor is mid-charge, for example, will observe
that capacitor's voltage reset rather than continue from its prior value. Carrying reactive
state across a topology rebuild in the same manner the memristor's state is carried is
planned future work and is primarily a bookkeeping problem: the new topology's node indices
must be matched to the element's previous node indices before its stored voltage or current
can be meaningfully re-applied.

\textbf{Fixed presets rather than continuous values.} Every adjustable quantity -- a
resistor's resistance, a source's amplitude or frequency -- is one of a small number of
discrete presets cycled by right-clicking, rather than a numeric entry. This design keeps
the interaction model consistent with the remainder of the base game (no custom text-entry
interface is required), at the cost of the fine-grained control an actual bench instrument
provides; a numeric-entry interface for these values is a natural extension.

\textbf{Solver scalability.} The MNA system for each independent network is solved by dense
Gaussian elimination every tick, which is adequate for the circuit sizes a single player
typically constructs by hand, but scales as $O(n^3)$ in the number of independent nodes and
does not exploit the sparsity that a large circuit's conductance matrix would in practice
possess. A sparse solver, or amortization of re-solves across ticks when no electrically
relevant change has occurred, would be necessary before the tool could comfortably support
very large student-built circuits, or a heavily populated multiplayer server hosting
numerous independent networks simultaneously.

\textbf{Verification remains informal.} As noted in Section~\ref{sec:validation}, the
closed-form comparisons used during development have not yet been captured as a persisted,
automatically re-run regression suite, and the tool's pedagogical effectiveness has to date
been assessed only informally by its author, rather than through a controlled classroom
study. A natural next step is a pre/post evaluation using an established circuits concept
inventory, comparing a cohort using CircuitCraft as a supplementary laboratory activity
against a cohort using a conventional schematic-based simulator, in order to test whether
the game-based, first-person construction experience described in Section~\ref{sec:pedagogy}
produces a measurable difference in conceptual understanding, rather than in engagement
alone.

\textbf{Memristor model fidelity.} As discussed in Section~\ref{sec:related-work}, the mod
implements only the original HP linear-drift formulation, which is known to exhibit boundary
pinning under sustained excitation, a limitation subsequent compact models such as Biolek et
al.'s windowed variant~\cite{biolek2009spice} were specifically designed to address. Exposing
one or more alternative compact models as an additional preset would allow the tool to also
illustrate why the modeling literature moved beyond the original formulation, rather than
presenting it as the final word on memristor behavior.

\textbf{Platform dependency.} The mod is necessarily tied to a specific version of
Minecraft and of the Fabric mod loader, both of which are externally maintained and continue
to evolve; Section~\ref{sec:validation} discusses the verification discipline this
requires, but ongoing maintenance to track future game updates remains an unavoidable cost
of building instructional software atop an actively developed commercial platform, rather
than a purpose-built, stable simulator.
